\documentclass[12pt]{article}
\usepackage[utf8]{inputenc}
\usepackage[T1]{fontenc}
\usepackage[margin=1in]{geometry}
\usepackage{setspace}
\usepackage{xcolor}
\usepackage[colorlinks=true, urlcolor=blue!60!black, linkcolor=blue!60!black]{hyperref}
\setlength{\parindent}{0em}
\setlength{\parskip}{0.45em}
\renewcommand{\baselinestretch}{1.0}
\begin{document}
\begin{center}
{\Large \textit{Pike v. Bruce Church, Inc.}, 397 U.S. 137 (1970)\par}
\vspace{0.3em}
{\normalsize Case Brief --- One-Page Edition\par}
\end{center}

\textbf{Memory Jogger.} The cantaloupe case and the \textit{Pike} balancing test for evenhanded burdens on interstate commerce (\href{https://supreme.justia.com/cases/federal/us/397/137/}{397 U.S. 137}).

\textbf{Facts.} Arizona ordered Bruce Church to pack its superior Parker cantaloupes in Arizona rather than at its sheds 31 miles away in California, which would have required a \$200,000 plant.

\textbf{Procedural History.} A three-judge district court enjoined the order; Arizona appealed directly.

\textbf{Votes.} Affirmed, 8--0: Justice Stewart for a unanimous Court.

\textbf{Holding.} An evenhanded law with incidental effects stands unless its burden is ``clearly excessive in relation to the putative local benefits''; Arizona's reputational interest could not justify forcing an in-state plant.

\textbf{Reasoning.} The order sought only ``reflected good will'' for other growers, an interest ``minimal, at best,'' and forcing operations in-state that could be done more efficiently elsewhere is ``virtually per se illegal.'' \textit{National Pork Producers Council v. Ross} (2023) later fractured over the test.

\textbf{Examples.} \emph{Same-Side:} in-state shrimp peeling mandates. \emph{Opposite-Side:} an evenhanded polystyrene packaging ban. \emph{Fence-Sitter:} a ban on milk from hormone-treated cows.

\textbf{Critique.} The decision rests more on the anti-local-processing rule than true balancing, and it offers no metric for weighing costs against noneconomic values.

\textbf{Key Quotation.} ``Although the criteria for determining the validity of state statutes affecting interstate commerce have been variously stated, the general rule that emerges can be phrased as follows: where the statute regulates evenhandedly to effectuate a legitimate local public interest, and its effects on interstate commerce are only incidental, it will be upheld unless the burden imposed on such commerce is clearly excessive in relation to the putative local benefits. \textit{Huron Cement Co. v. Detroit}, 362 U. S. 440, 443. If a legitimate local purpose is found, then the question becomes one of degree. And the extent of the burden that will be tolerated will, of course, depend on the nature of the local interest involved, and on whether it could be promoted as well with a lesser impact on interstate activities.'' --- Justice Stewart, for the Court, 397 U.S. at 142.

\end{document}
